\documentclass[11pt,]{article}
\usepackage{lmodern}
\usepackage{amssymb,amsmath}
\usepackage{ifxetex,ifluatex}
\usepackage{fixltx2e} % provides \textsubscript
\ifnum 0\ifxetex 1\fi\ifluatex 1\fi=0 % if pdftex
  \usepackage[T1]{fontenc}
  \usepackage[utf8]{inputenc}
\else % if luatex or xelatex
  \ifxetex
    \usepackage{mathspec}
  \else
    \usepackage{fontspec}
  \fi
  \defaultfontfeatures{Ligatures=TeX,Scale=MatchLowercase}
    \setmainfont[]{Fira Sans}
\fi
% use upquote if available, for straight quotes in verbatim environments
\IfFileExists{upquote.sty}{\usepackage{upquote}}{}
% use microtype if available
\IfFileExists{microtype.sty}{%
\usepackage{microtype}
\UseMicrotypeSet[protrusion]{basicmath} % disable protrusion for tt fonts
}{}
\usepackage[margin=.5in]{geometry}
\usepackage{hyperref}
\hypersetup{unicode=true,
            pdfborder={0 0 0},
            breaklinks=true}
\urlstyle{same}  % don't use monospace font for urls
\usepackage{natbib}
\bibliographystyle{plainnat}
\usepackage{longtable,booktabs}
\IfFileExists{parskip.sty}{%
\usepackage{parskip}
}{% else
\setlength{\parindent}{0pt}
\setlength{\parskip}{6pt plus 2pt minus 1pt}
}
\setlength{\emergencystretch}{3em}  % prevent overfull lines
\providecommand{\tightlist}{%
  \setlength{\itemsep}{0pt}\setlength{\parskip}{0pt}}
\setcounter{secnumdepth}{0}

%%% Use protect on footnotes to avoid problems with footnotes in titles
\let\rmarkdownfootnote\footnote%
\def\footnote{\protect\rmarkdownfootnote}


  \title{}
    \author{}
    \date{}
  

%%%%%%%%%%
% personal preamble edits here
%%%%%%%%%%
\pagenumbering{gobble}

% can toggle this for Helvetica
%\usepackage{helvet}
%\renewcommand{\familydefault}{\sfdefault}

% \titlespacing*{\paragraph}{0pt}{2pt}{1em}

% set section numbering/lettering
% tips for seccntformat: https://tex.stackexchange.com/questions/95896/how-to-format-subsection-title-without-packages
\makeatletter
\def\@seccntformat#1{%
  \expandafter\ifx\csname c@#1\endcsname\c@section\else
  \expandafter\ifx\csname c@#1\endcsname\c@paragraph\else
  \csname the#1\endcsname\quad
  \fi\fi}
  
  % stucture for these commands: https://texfaq.org/FAQ-atsigns
  \renewcommand\section{
  \@startsection{section}{1}{\z@}
    {-3.5ex \@plus -1ex \@minus -.2ex}
    {1.0ex \@plus.2ex} %reduce space below section (was 1.5ex)
    {\normalfont\normalsize\bf\uppercase}} %modify font style
    
  \renewcommand\subsection{
  \@startsection{subsection}{2}{\z@}
    {-1.5ex\@plus -1ex \@minus -.2ex}%reduce space above subsection (was -3.25ex)
    {0.5ex \@plus .2ex}%reduce space below subsection (was 1.5ex)
    {\normalfont\normalsize\bf}} %modify font style
    
  \renewcommand\subsubsection{
  \@startsection{subsubsection}{3}{\z@}
    {-1.0ex\@plus -1ex \@minus -.2ex}%reduce space above subsubsection (was -3.25ex)
    {0.5ex \@plus .2ex}%reduce space below subsubsection (was 1.5ex)
    {\normalfont\normalsize\bf}} %modify font style
    
  \renewcommand\paragraph{
  \@startsection{paragraph}{4}{\z@}
    {-0.5ex\@plus -1ex \@minus -.2ex}%reduce space above paragraph (was -3.25ex)
    {-1.5ex \@plus .2ex}%convert space below paragraph to an indent (was 1.5ex)
    {\normalfont\normalsize\bf}} %modify font style    
\makeatother

\renewcommand\thesubsection{\Alph{subsection}.}
\renewcommand\thesubsubsection{\thesubsection\arabic{subsubsection}.}

% reduce spacing at the top of lists
\usepackage{enumitem}
\setlist{topsep = 2pt}

% allow text to wrap around figures
\usepackage{graphicx}
\usepackage{wrapfig}

%%%%%%%%%%

\begin{document}

\hypertarget{nimh-grant-1-r03-mh123787-01}{%
\paragraph*{\texorpdfstring{\textbf{NIMH Grant}: 1 R03
MH123787-01}{NIMH Grant: 1 R03 MH123787-01}}\label{nimh-grant-1-r03-mh123787-01}}
\addcontentsline{toc}{paragraph}{\textbf{NIMH Grant}: 1 R03 MH123787-01}

\hypertarget{title-gwas-of-the-rdoc-cognitive-systems-domain-modeling-the-latent-genetic-architecture-of-working-memory}{%
\paragraph*{\texorpdfstring{\textbf{Title}: GWAS of the RDoC Cognitive
Systems Domain: Modeling the Latent Genetic Architecture of Working
Memory}{Title: GWAS of the RDoC Cognitive Systems Domain: Modeling the Latent Genetic Architecture of Working Memory}}\label{title-gwas-of-the-rdoc-cognitive-systems-domain-modeling-the-latent-genetic-architecture-of-working-memory}}
\addcontentsline{toc}{paragraph}{\textbf{Title}: GWAS of the RDoC
Cognitive Systems Domain: Modeling the Latent Genetic Architecture of
Working Memory}

\hypertarget{pi-joey-trampush}{%
\paragraph*{\texorpdfstring{\textbf{PI:} Joey
Trampush}{PI: Joey Trampush}}\label{pi-joey-trampush}}
\addcontentsline{toc}{paragraph}{\textbf{PI:} Joey Trampush}

\hypertarget{project-period-06152020-06142022}{%
\paragraph*{\texorpdfstring{\textbf{Project Period}: 06/15/2020 --
06/14/2022}{Project Period: 06/15/2020 -- 06/14/2022}}\label{project-period-06152020-06142022}}
\addcontentsline{toc}{paragraph}{\textbf{Project Period}: 06/15/2020 --
06/14/2022}

\hfill\break
\hfill\break
This grant was funded in response to NIMH PAR-17-158, ``Secondary Data
Analyses to Explore NIMH Research Domain Criteria RDoC (R03).'' This
report specifies activities from June 15, 2020 through June 14, 2021.

\hypertarget{b.-accomplishments}{%
\section{B. Accomplishments}\label{b.-accomplishments}}

\hypertarget{b.1-what-are-the-major-goals-of-the-project}{%
\subsection*{B.1 What are the major goals of the
project?}\label{b.1-what-are-the-major-goals-of-the-project}}
\addcontentsline{toc}{subsection}{B.1 What are the major goals of the
project?}

\textbf{Working memory} is the cognitive system responsible for the
transient holding and processing of information in conscious awareness
just long enough to use it for problem solving. Limited working memory
faculties are a fundamental aspect of the cognitive impairments
prevalent in ADHD, schizophrenia, and many other CNS disorders. More
than half of the variability in working memory capacity from person to
person can be sourced back to inherited genetic factors. Thus, our goal
is to fast-track the identification and localization of molecular
genetic signatures of working memory, and then to uncover the underlying
causal molecular and biological mechanisms linking working memory
deficits to genetically correlated disorders. GWAS (genome-wide
association study) is a powerful approach to identify molecular
signatures of complex traits if adequate sample sizes and proper
statistical procedures are employed. We will utilize existing genomic
and phenotypic data from the Cognitive Genomics Consortium (COGENT) to
test novel hypotheses about the molecular genetic structure of working
memory within the flexible RDoC framework. The study PI is an
investigator and data analyst for COGENT, which was formed to conduct
large-scale GWAS of cognitive traits implicated in schizophrenia such as
general cognitive ability, processing speed, working memory, and
episodic memory. At present, we have individual-level GWAS data on
24,000 participants who have 100,000 working memory data points
collected for analysis, which will be the largest and most well powered
GWAS of working memory to date. The following interrelated specific aims
are proposed:

\textbf{Aim 1. Conduct a common factor GWAS of working memory}.

\emph{Method}. Genomic structural equation modeling (Genomic SEM) will
be used to ascertain the joint genetic architecture of narrow working
memory facets to identify loci with effects on a genetically-derived
general latent working memory factor (``\emph{Gwm}'').
\emph{Hypothesis}. With sample sizes in the tens of thousands, the first
replicable GWAS hits for working memory are expected to emerge.

\textbf{Aim 2. Determine the extent of shared genetic architecture}.

\emph{Method}. At genome-wide scale, linkage disequilibrium score
regression (LDSC) and genomic restricted maximum likelihood (GREML)
estimation will be used to determine the range of pleiotropy between
working memory and correlated CNS phenotypes. \emph{Hypothesis}.
Widespread coheritability is expected to emerge, and \emph{Gwm} will
likely outperform narrower working memory domains and individual working
memory subtests.

\textbf{Aim 3: Isolate the strongest causal loci}.

\emph{Method}. At the molecular level, Mendelian randomization (MR) will
be used to determine the direction of causality underlying pleiotropy
between working memory and correlated CNS phenotypes. \emph{Hypothesis}.
\emph{Gwm} hits are expected to exhibit strong causal effects from
working memory to ADHD and schizophrenia, and moderate causal effects to
other CNS phenotypes.

\textbf{Aim 4. Map key causal biological mechanisms}.

\emph{Method}. To prioritize \emph{Gwm} loci for follow-up studies,
functional mapping and annotation of genetic associations (FUMA) will be
used to characterize the biological function of causal loci.
\emph{Hypothesis}. Elucidating the biology of working memory is expected
to be streamlined by narrowing the search space to causal CNS
mechanisms.

\hypertarget{b.1.a-have-the-major-goals-changed-since-the-initial-competing-award-or-previous-report}{%
\subsubsection{B.1.a Have the major goals changed since the initial
competing award or previous
report?}\label{b.1.a-have-the-major-goals-changed-since-the-initial-competing-award-or-previous-report}}

No

\hypertarget{b.2-what-was-accomplished-under-these-goals}{%
\subsection{B.2 What was accomplished under these
goals?}\label{b.2-what-was-accomplished-under-these-goals}}

\hypertarget{challenges}{%
\subsubsection{Challenges}\label{challenges}}

Year 1 was challenged by the pandemic and slowed our launch. We're
catching up now, but had restricted access to our labs at USC for almost
the entire year and limited access to people/colleagues/students. We
rode it out and thankfully are back on track, and are looking forward to
generating the results as promised in Year 2.

Also back in August, the USC Center for Advanced Research Computing
(CARC) launched a new high-performance computing cluster, Discovery,
which is what we use to compute our large-scale genomic analyses. The
new cluster now consists of around 11,000 CPU cores in around 500
compute nodes; the typical compute node has dual 8 to 16 core processors
and resides on a 56 gigabit FDR InfiniBand backbone; and nearly 100 GB
of storage space is provided for each user, with a default storage quota
of 10 TB per PI across our projects. However, as expected, the new
cluster has had its share of ups and downs in terms of operating
smoothly and wihtout interruption (e.g., it has been offline for
maintenance for several weeks at a time here and there this year, just
to work out some of the unforseeable bugs). There was also a bit of a
learning curve to orient to the new user interfaces and job submission
protocols. To support advanced computing at the university, USC recently
upgraded the network capacity of individual research buildings to 100
Gbps. The CARC has also deployed multiple 100 Gbps links between the ITS
data center and the rest of the campus, creating an aggregate bandwidth
of up to 400 Gbps. CARC will ultimately be an invaluable resource for
computational researchers at USC, but the newness of it just added some
additional wait time for current project.

Finally, The Trans-Omics for Precision Medicine (TOPMed) programme seeks
to elucidate the genetic architecture and biology of heart, lung, blood
and sleep disorders, with the ultimate goal of improving diagnosis,
treatment and prevention of these diseases. The initial phases of the
programme focused on whole-genome sequencing of individuals with rich
phenotypic data and diverse backgrounds. Here we describe the TOPMed
goals and design as well as the available resources and early insights
obtained from the sequence data. The resources include a variant
browser, a genotype imputation server, and genomic and phenotypic data
that are available through dbGaP (Database of Genotypes and
Phenotypes)1. In the first 53,831 TOPMed samples, we detected more than
400 million single-nucleotide and insertion or deletion variants after
alignment with the reference genome. Additional previously undescribed
variants were detected through assembly of unmapped reads and customized
analysis in highly variable loci. Among the more than 400 million
detected variants, 97\% have frequencies of less than 1\% and 46\% are
singletons that are present in only one individual (53\% among unrelated
individuals). These rare variants provide insights into mutational
processes and recent human evolutionary history. The extensive catalogue
of genetic variation in TOPMed studies provides unique opportunities for
exploring the contributions of rare and noncoding sequence variants to
phenotypic variation. Furthermore, combining TOPMed haplotypes with
modern imputation methods improves the power and reach of genome-wide
association studies to include variants down to a frequency of
approximately 0.01\%.

\begin{longtable}[]{@{}llll@{}}
\toprule
Variable & TOPMed (v.R2) & HRC (v.r1.1) & 1000G Phase 3 (v.5) \\
\midrule
\endhead
Haplotypes & 194,512 & 64,940 & 5,008 \\
Samples & 97,256 & 32,470 & 2,504 \\
Sites (chr1-22) & 308,107,085 & 39,635,008 & 49,143,605 \\
\bottomrule
\end{longtable}

\hypertarget{opportunities}{%
\subsubsection{Opportunities}\label{opportunities}}

\hypertarget{dockersingularity}{%
\paragraph{Docker/Singularity}\label{dockersingularity}}

\hypertarget{data-visualizationthe-penn-letter-n-back-test}{%
\paragraph{Data visualization/The Penn Letter N-Back
Test}\label{data-visualizationthe-penn-letter-n-back-test}}

The Penn Letter N-Back Test measures working memory (WM), the ability to
keep and refresh goal-related information. Participants attend to a
continual series of letters that flash on the screen (one at a time) and
press the spacebar according to three different rules (called the
0-back, 1-back, and 2-back). During the 0-back condition, the
participant must simply respond to a currently present target (``X'').
During the 1-back condition, he/she must press the spacebar when the
letter on the screen is the same as the previous letter. During the
2-back condition, he/she must press the spacebar when the letter on the
screen is the same as the letter before the previous letter (i.e.~2
letters back). In all trials, the inter-stimulus interval (ISI) is 2.5
seconds, and the stimuli (letters) themselves are presented for 0.5
seconds each. The participant practices all three principles before
testing.

\hypertarget{b.3-competitive-revisionsadministrative-supplements}{%
\subsection{B.3 Competitive Revisions/Administrative
Supplements}\label{b.3-competitive-revisionsadministrative-supplements}}

\hypertarget{for-this-reporting-period-is-there-one-or-more-revisionsupplement-associated-with-this-award-for-which-reporting-is-required}{%
\subsubsection{For this reporting period, is there one or more
Revision/Supplement associated with this award for which reporting is
required?}\label{for-this-reporting-period-is-there-one-or-more-revisionsupplement-associated-with-this-award-for-which-reporting-is-required}}

No

\hypertarget{b.4-what-opportunities-for-training-and-professional-development-has-the-project-provided}{%
\subsection{B.4 What opportunities for training and professional
development has the project
provided?}\label{b.4-what-opportunities-for-training-and-professional-development-has-the-project-provided}}

Many many. Reproducible research. R. Rmarkdown. GitHub. Git.

Support and Training User support for advanced computing at USC is
primarily available through the CARC, which consists of a systems team,
cloud/DevOps team, and research support team. The staff has expertise in
storage system administration, software development, science and
engineering applications, high-performance and parallel programming, and
cloud computing solution development. The CARC also works with
researchers who do not require HPC-level performance but need help
developing or optimizing their code, improving their computing
productivity by learning better programming languages, scripting, and so
forth. The CARC offers regular onboarding sessions, office hours,
introductory and advanced training workshops, and in-person and remote
consultations. It also hosts vendor and national supercomputing center
training at USC. As XSEDE Campus Champions, CARC staff can connect
researchers with high-end computing needs to national supercomputing
centers to scale their workflows and reach other researchers nationwide.

\hypertarget{b.5-how-have-the-results-been-disseminated-to-communities-of-interest}{%
\subsection{B.5 How have the results been disseminated to communities of
interest?}\label{b.5-how-have-the-results-been-disseminated-to-communities-of-interest}}

Neuropsychopharmacology

\hypertarget{b.6-what-do-you-plan-to-do-during-the-next-reporting-period-to-accomplish-the-goals}{%
\subsection{B.6 What do you plan to do during the next reporting period
to accomplish the
goals?}\label{b.6-what-do-you-plan-to-do-during-the-next-reporting-period-to-accomplish-the-goals}}

\begin{wrapfigure}{R}{0.3\textwidth}

\hfill{}\includegraphics[width=.3\textwidth]{rppr_trampush_nimh_grant_R03MH123787_files/figure-latex/unnamed-chunk-2-1} 

\caption{Important scatterplot}\label{fig:unnamed-chunk-2}
\end{wrapfigure}

\hypertarget{timeline}{%
\subsubsection{Timeline}\label{timeline}}

There are lots of good examples of \texttt{R}-based Gantt charts to be
found by clever Googling. For displaying progress with sidebar
annotations by aim, I particularly like
\href{https://datascienceplus.com/visualize-your-cvs-timeline-with-r-gantt-style/}{\underline{this}}
example from the
\href{https://github.com/laresbernardo/lares}{\underline{lares}}
package.

\bibliography{references.bib}


\end{document}
